\documentclass{slides}

\title{Relation algebra and finite axiomatization of NF(U)}

\author{Lavinia Randall Holmes\\Boise State University (retired)}

\date{June 9, 2026\\ at the Edmonton Logic Seminar}

\begin{document}

\begin{slide}

\maketitle

\end{slide}

\begin{slide}

That stratified comprehension can be finitely axiomatized is an old result, due to Theodore Hailperin.  His construction is analogous to the finite axiomatization of predicate class theory.  His axiomatization is notably awful.

Other finite axiomatizations have been presented since.

The one described here was used in the first chapters of my book Elementary Set Theory with a Universal Set.  It is adapted from the axiomatization of relation algebra combined with the interpretation of first order logic in relation algebra with ordered pair relations, due to Tarski and Givant and previously discussed in this seminar.

\end{slide}

\begin{slide}


The primitive notions of the theory we describe are membership, equality,  the projection relations $\pi_1$ and $\pi_2$ of the ordered pair [taken as a primitive notion] and the empty set $\emptyset$.

The first axiom, which plays no essential role in our development, is weak extensionality.  Strong extensionality could be used if desired, and one would obtain NF instead of NFU.

The weak extensionality axiom asserts that nonempty sets are determined exactly by their extensions.

$$(\forall ABz:z\in A \rightarrow (A=B \leftrightarrow (\forall x:x\in A \leftrightarrow x \in B)))$$

\end{slide}

\begin{slide}

We now begin the axiomatization of stratified comprehension.  The intention is that the axioms we present seem
simply to be natural hypotheses about what sets exist.

\begin{description}

\item[Axiom of the Empty Set:]  We provide a primitive constant $\emptyset$ with axiom $(\forall x:\neg x \in \emptyset)$.

\item[Definition:]  We say that $x$ is a set iff $$(\exists y:y \in x) \vee x = \emptyset.$$   We can then define
$\{x\mid \phi\}$ as the unique set $A$ such that $(\forall x \mid x \in A \leftrightarrow \phi)$, for any formula $\phi$.  This is an inessential but very useful strengthening of the original conceptual kit of NFU.

\end{description}

\end{slide}

\begin{slide}

\begin{description}

\item[Axiom of Singletons:]  For every $x$, $$\{x\} = \{y\mid y=x\}$$ exists.    We also use the notation $\iota(x)$ for $\{x\}$, which has the advantage that $\iota^n(x)$ can be used to represent repeated application.

To an audience of philosophers, I should emphasize that no existence predicate is mentioned here.  ``$\{x \mid \phi\}$ exists'' is a convenient way to
say $$(\exists A:(\forall x:x \in A \leftrightarrow \phi))$$ ($A$ not free in $\phi$).

It is worth noticing that the axiom of singletons and the axiom of the empty set are together enough to show that there are infinitely many objects in any model of NFU [the fact which Quine persisted in saying is evidence for the axiom of infinity in NF, which it does not, as Infinity is not provable in NFU].

\end{description}

\end{slide}

\begin{slide}

\begin{description}


\item[Axiom of Complements:]  For each set $A$, $A^c = \{x\mid x \not\in A\}$ exists.

\item[Axiom of Binary Union:]  For any sets $A,B$, $A \cup B = \{x \mid x \in A \vee x \in B\}$ exists.

\item[Definition:]  $V$ is defined as $\emptyset^c = \{x\mid x = x\}$.

\item[Definition:]  $A \cap B$ is defined as $(A^c \cup B^c)^c$ and is equal to $\{x\mid x \in A \wedge x \in B\}$.
$A-B$ is defined as $A \cap B^c = \{x \in A \wedge\neg A \in B$.

\item[Definition:]  $\{x_1,x_2\}$ is defined as $\{x_1\} \cup \{x_2\}$,  $\{x_1,x_2,\ldots,x_n\}$ is defined
as $\{x_1\} \cup \{x_2,\ldots,x_n\}$.

This gives us the Boolean algebra of all concrete finite sets (of whatever objects there are) and their complements.

\end{description}

\end{slide}

\begin{slide}

We did mention relations.

\begin{description}

\item[Axiom of Ordered Pairs:]  For each $x,y$ there is a unique $z$ such that $z \pi_1 x$ and $z \pi_2 y$, which we
call $(x,y)$.  For all $x,y,z,w$, $(x,y)=(z,w)$ implies $x=z$ and $y=w$. Objects $(x,y)$ are called ordered pairs.

\item[Definition:]  A relation is a set of ordered pairs.  We write $\{(x,y)\mid\phi\}$ for $$\{z\mid (\exists xy:z=(x,y) \wedge \phi)\},$$ where $z$ is a fresh variable.

We write $x \,R\,y$ for $(x,y) \in R$, when $R$ is a relation.

We comment that in our book we introduced ordered pairs as a primitive in the context of NFU [the intention being that the pair be type level, which cannot be achieved by definition in NFU].

\end{description}


\end{slide}
\begin{slide}

\begin{description}

\item[Axiom of Cartesian Products:]  For any sets $A,B$, $A \times B =\{(a,b):a \in A \wedge b \in B\}$ exists.  We define $V^2$ as $V \times V$ and $V^{n+1}$ as $V \times V^n$.

\item[Definition:]  When $R$ is a relation, we write $R^c$ for $(V\times V)-R$.  The overloading should be harmless
as long as we are clear about the intentions with which letters are used.

\item[Axiom of Converses:]  For any relation $R$, $R^{-1} = \{(y,x)\mid (x,y)\in R\}$ exists.

\item[Axiom of Relative Products:]  For any relations $R,S$, $R|S = \{(x,y)\mid (\exists z:x\,R\,z \wedge z\,S\,y)\}$ exists.

\item[Axiom of the Diagonal:]  $\{(x,x)\mid x=x\}$ exists.

\end{description}

\end{slide}

\begin{slide}

At this point we have defined the primitives of relation algebra.  To support the strategy of Tarski and
Givant, we need the projection relations.  To support set definitions involving membership, we need the relation of inclusion. 

\begin{description}

\item[Axiom of Projections:]  $[\pi_1]=\{(x,y)\mid x \pi_1 y\}$ exists.   $[\pi_2]=\{(x,y)\mid x \pi_2y\}$ exists.  I have always said that Ordered Pairs implies Infinity, but in this axiomatization it is actually this axiom that entails Infinity.

\item[Axiom of Inclusion:]  $$[\subseteq] = \{(x,y):(\forall z:z \in x \rightarrow z \in y\}$$ exists.

\end{description}

It should be noted, though we will not appeal to this, that we now have everything needed for the interpretation
of first order logic in relation algebra due to Tarski and Givant.

\end{slide}

\begin{slide}

We need three additional axioms, one to translate from the domain of relations to the domain of sets,
and one to establish parallelism between the domain of general objects and the domain of singleton sets.  The last is the familiar axiom of set union.

\begin{description}

\item[Axiom of Domains:]  For any relation $R$, ${\tt dom}(R) = \{x\mid (\exists y:x \,R\,y)\}$ exists.

\item[Axiom of Singleton Images:]  For any relation $R$, the relation $R^\iota = \{\{x\},\{y\}):x\,R\,y\}$ exists.  Note that ${\tt dom}((A \times V)^{\iota})$ implements $\iota``A = \{\{x\}\mid x \in A\}$.

\item[Axiom of Set Union:]  For any set $A$, $\bigcup A = \{x\mid(\exists a \in A:x \in a \wedge a \in A)\}$ exists.  If set union is omitted, we get predicative NFU.

\end{description}

\end{slide}

\begin{slide}

This completes our list of axioms.   We introduced these as the axioms of the theory in our book for rhetorical purposes.  These are all quite natural hypotheses about the existence of sets and relations, and not only believable but motivated by the need to reason about relations.  We do not give stratified comprehension (a quite esoteric scheme) as the motivation of these axioms.  But we proceed to prove that the scheme of stratified comprehension follows from the axioms, and so is supported by any intuition which finds the axioms reasonable

\end{slide}

\begin{slide}

In what follows, we will show that the scheme of stratified comprehension follows from our axioms.  In the process,
we will in effect establish that first order logic can be interpreted via relation algebra with projection relations, though we are not exactly giving a proof of the results of Tarski and Givant, because we allow ourselves to freely use the interplay between sets and relations in our axiomatics.

\end{slide}

\begin{slide}

We say that a sentence $P(x_1,\ldots,x_n)$ with the $x_i$'s all its free variables is interpreted by a set
$A$ just in case $A=\{(x_1,\ldots,x_n):P(x_1,\ldots,x_n)\}$.  We stipulate that we define $(x_1,x_2,\ldots,x_n)$
as $(x_1,(x_2\ldots,x_n))$.  We could conveniently define $(x_1)$ as $x_1$ and () as a convenient fixed object such as the empty set.

Some formulas (for example $x_1 \not\in x_1$) are well known not to be interpreted by any set.

\end{slide}

\begin{slide}

If $P(x_1,\ldots,x_n)$ and $Q(x_1,\ldots,x_n)$ have the same free variables (which we leave out of the notation for the rest of this slide)  then $\neg P$ is interpreted by $\{(x_1,\ldots,x_n)\mid P\}^c$ and $P \vee Q$ is interpreted
by $\{(x_1,\ldots,x_n)\mid P\} \cup \{(x_1,\ldots,x_n)\mid Q\}$.  This generalizes to all operations of propositional logic.

\end{slide}

\begin{slide}

If $P(x_1,\ldots,x_n)$ is interpreted, then $(\exists x_1: P)$ is interpreted by ${\tt dom}(\{(x_1,\ldots,x_n)\mid P\}^{-1})$, which has free variables $(x_2,\ldots,x_n)$.  The universal quantifier is of course definable in terms of the existential quantifier and negation.

\end{slide}


\begin{slide}

For any relation $R$, we define $R^0$ as $[=]$ and $R^{n+1}$ as $R^n|R$.

We define $\Pi_n$ as $[\pi_2]^{n-1}|[\pi_1]$.  $(x_1,\ldots,x_i,\ldots x_n)$ has the relation $\Pi_i$ to each $x_i$ for $i<n$
and the relation $\Pi_n^* = [\pi_2]^{n-1}$ to $x_n$.  We define $\Pi^m_n$ ($m \geq n$) as $\Pi_n$ if $m>n$ and $\Pi^*_n$ if $m=n$.

We can then define the relation holding between a tuple $(x_1,\ldots,x_n)$ and a tuple $(y_1,\dots,y_m)$
where $x_i = y_{f(i)}$ for each $i$ as the intersection of all relations $\Pi^n_i|(\Pi^m_{f(i)})^{-1}$, for $1 \leq i \leq n$, starring the $\Pi$ operators where appropriate. 

We can represent the image $R``A$ of a set $A$ under a relation $R$ as ${\tt dom}((R \cap (A\times V))^{-1})$.

\end{slide}

\begin{slide}

We can then represent $R(y_{f(1)},\ldots,y_{f(n})$ as a predicate of $(y_1,\ldots,y_m)$, as $$[\bigcap_{1 \leq i\leq n}(\Pi^n_i|(\Pi^m_{f(i)})^{-1})]``R \cap V^m).$$
 Note that this readily allows us to represent quantification over variables other than $x_1$, by reordering first and then undoing the reordering if desired).

In this way we can make every constituent formula of a complex formula be represented have the same argument list (it is also convenient to arrange for bound variables to be identified only with variables bound by the same binder).  We can use the same construction to remove variables which in fact have no effect.

This has the effect that representable formulas are closed under all constructions of first order logic.

We note that the membership relation $[\in]$ is not representable:  if it
were, we could construct ${\tt dom}([\in] \cap [=])^c$, and enough said.
\end{slide}


\begin{slide}

We transform any stratifiable formula $\phi$ with type assignment $\sigma$ into a representable formula, illustrating the use of the final elements of our toolkit.

Recall that we introduced the notation $\iota(x)$ for $\{x\}$ because we can conveniently write $\iota^n(x)$ for iteration of this operation $n$ times.

Where $N$ is the maximum element of the range of $\sigma$, replace each variable $x$  in the formula $\phi$ with a variable $x^*$ representing $\iota^{N-\sigma(x)}(x)$.  Variables continue to have the same reference whereever they appear, because they always appear at the same type.  A bound variable $x^*$ is of course supposed restricted to $\iota^{N-\sigma(x)}``V$ (add an atomic formula
$x^* \in \iota^{N-\sigma(x)}``V$ connected appropriately to restrict the quantifier).

To preserve the meaning of the formula, we need to replace the relations in atomic formulas.
\end{slide}

\begin{slide}

Any atomic formula $x \in y$ is then interpreted as $$\iota^{N-\sigma(x)}(x) ???\iota^{N-\sigma(y)}(y),$$  where ??? is a relation to be determined.  Note that this
is  $\iota^{N-\sigma(x)}(x) ??? \iota^{N-(\sigma(x)+1)}(y)$ and can be helpfully written $\iota^M(x) ??? \iota^{M-1}(y)$ (where $M=N-\sigma(x)$)  Now if $M=1$ (the smallest possible value) we want $\{x\} ??? y$ to be equivalent to $x \in y$ in the original sense.
The relation to use is the inclusion relation with domain restricted to the collection of singletons.  The collection of singletons
is $\iota``V= {\tt dom}(V \times V)^{\iota}$, so the desired relation of ``atomic inclusion'' is $E = (\iota``V \times V) \cap [\subseteq]$.
Then ??? in the general formula $x\in y$ read as $\iota^M(x) ??? \iota^{M-1}(y)$ is $E^{\iota^{M-1}}$ (the notation indicating iteration of the singleton image operation on relations $M-1$ times).  

\end{slide}

\begin{slide}
Other atomic relations are simply replaced by suitable iterated singleton images:  the other relations are easier because they are representable and will be flanked by variables assigned the same type, so the same number of iterated applications of the singleton operation.

\end{slide}
\begin{slide}
This immediately shows that stratified comprehension holds.  A stratified formula $\phi(x)$ converts to a representable formula
$\phi(x^*)$ where $x^*$ represents $\iota^M(x)$ for some $M$.  From its representation, it is straightforward to construct $\{\iota^M(x)\mid \phi\}$ and then application of Set Union
$M$ times gives $\{x\mid \phi\}$.

Someone might think of parameters in $\phi$ here.  We note that $x=a$ for any free variable $a$ is represented by $\{(x,a)\mid x=x\}$, which exists by our axioms, being $V \times \{a\}$.  This applies in particular to the handling of the bounding of quantifiers to iterated singleton images of the universe during the formula transformation.
\end{slide}

\begin{slide}
We do admit that we have given no more than a hint of the Tarski and Givant result, because we actually make free use of the interplay
between sets and relations afforded by the domain operation.  But I suspect it is a quite strong hint as to how their result is proved.

We think that we have entirely dispelled a common criticism of NF(U).  Stratified comprehension is not simply a syntactical trick.  Here we see that it is a consequence of a finite list of principles for constructing sets which are quite credible, though not the same as those in the Zermelo toolkit we are used to.

\end{slide}
\begin{slide}

A technical footnote:  we axiomatized NFU with a type level ordered pair, which is an inessential strengthening of NFU + Infinity.  NFU without Infinity can be axiomatized along these lines.  Let $(x,y)$ represent the usual Kuratowski pair
$\{\{x\},\{x,y\}\}$.  The projections of the Kuratowski pair are provably not representable relations.
So restate the axiom of projections to provide $[\pi_i]$ as the set $$\{(x,y)\mid y = \iota^2(\pi_i(x)\}$$ for $i=1,2$.  Everything to do with projections and pairing is actually handled in a stratified formula $y$ by atomic formulas $x \,\pi_i\,y$ which can conviently be read $x [\pi_i] \iota^2(y)$ and then transformed into forms involving iterated singleton image operations as discussed above.  But we do not give full details of this adaptation, as it also involves redefining general $n$-tuples in a really horrible way to get all their components to be of the same relative type.   NFU with the type level ordered pair is a better foundational theory, anyway.

\end{slide}

\begin{slide}
If we adopt strong extensionality we have an axiomatization of NF.  NF proves that there is a type level ordered pair:  I think it is most efficient to use the primitive one and identify it with the Quine pair after the fact.  The reason for this is that the Quine pair definition is quite complicated, and we do not want to try to smuggle it in as a convincing basic axiomatic construction.

\end{slide}

\begin{slide}

We outline a purely set theoretic solution which works in NFU without Infinity (or in NF, perfectly well).  This uses not relation algebra
but a more direct representation of constructions of first order logic.

Initially, we adopt from our axiom list above the axioms of weak extensionality (strong extensionality would give NF), empty set, singletons, complement and binary union.

As a term of art for this construction, we call a set a {\em variable\/} iff it has more than one element.

An environment $e$  is a set of pairs $\{\{x\},v\}$ where $v$ is a variable, with the further property that for any variable $w$ there
is at most one $x$ such that $\{\{x\},w\}\in e$.  Where $e$ is an environment and $\{\{x\},w\}\in e$ we write $x=e[w]$.  Notice that
the empty set is a domain.

We provide as an axiom that for each environment $e$, the set ${\tt dom}(e)$ of variables which belong to elements of $e$ exists.
This is convenient but perhaps not necessary:  what we really need in what follows is a convenient way to say that two environments
have the same domain, and equality of these domain sets is a very handy way to say this.

\end{slide}

\begin{slide}

For any environments $e$ and $f$, we define $e+f$ as $e\cup f$ just in case it is an environment, and otherwise leave it undefined.
We do not need any additional axioms to show that $e+f$ exists (when it does).

A predicate is a set of environments all of which have the same domain, which we call the Domain of $P$, with the special proviso
that the empty set has every $d$ as its Domain.

For any predicate $P$ with Domain $d$, we provide by axiom the existence of the predicate $\neg P$ which contains all environments
with Domain $d$ which do not belong to $P$.

For any predicate $P$ with Domain $d$ and predicate $Q$ with Domain $f$, we provide a predicate $P \wedge Q$ which consists of
all environments $e+f$ with $e \in P$ and $f \in Q$.

Our choice of notation for these operations is not accidental.

\end{slide}

\begin{slide}

For any predicate $P$ with Domain $d$ and variable $v$, the predicate $(\exists v:P)$ is defined as $P$ if $v$ is not in the Domain of $P$,
and otherwise as the set of all environments $e$ with $v \not\in {\tt dom}(e)$ such that there is $x$ such that $e \cup \{\{\{x\},v\}\} \in P$.

We provide for any variables $v$ and $w$ the predicate $=_{vw} = \{\{\{\{x\},v\},\{\{x\},w\}\}\mid x=x\}$, the equality relation $v = w$.

We provide for any variables $v$ and $w$ the atomic inclusion relation $E_{v,w} = \{\{\{\{\{x\}\},v\},\{\{y\},w\}\}\mid x\in y\}$, representing
$v \in \iota``V \wedge v \subseteq w$.

\end{slide}

\begin{slide}

We provide the singleton image construction:  if $e$ is an environment, the set $e^\iota$ exists which contains $\{\{\{x\}\},\{v,\emptyset\}\}$
exactly if $e$ contains $\{\{x\},v\}$ [we have to do a little trick here with the variable to preserve stratification].  $P^\iota$ is defined as the set of $e^\iota$ such that $e \in P$, which is stipulated to exist by axiom.  The fact that we only use variables of special form in $P^\iota$ is irrelevant, as we can use quantification and equality to change the variables to whatever we like.

We provide that if $P$ is a predicate with Domain $\{v\}$, that $\{x\mid \{\{x\},v\} \in P\}$ exists (we provide the ability to recover
a set from a 1-ary predicate).

We provide the axiom of set union.

\end{slide}

\begin{slide}

It is evident from the way we define operations on predicates that we can represent any formula of first order logic in which each
atomic formula is represented by a predicate in an obvious sense by a set.

Thus, if we carry out the transformation on a stratified formula described above, amounting to replacing each variable with a suitable iterated singleton to bring everything to the same relative type, we can see that some iterated singleton image of any set defined by a stratified formula exists, and so by union that set exists.

Note that  restriction of values of a variable to $n$-ary singletons for any given $n$ is straightforward, and further that fixing the
value of a variable to any desired constant value is straightforward (this handles parameters).

\end{slide}

\begin{slide}

We think our approach here may have more general applicability.  Rather than building some ad hoc seeming toolkit of set theoretical operations which happens to handle representation of formulas, we provide a quite direct set theoretical representation of what it is
to be a predicate of a collection of free variables (a collection of assignments of values to the variables which make the predicate true)
and then provide the logical operations on predicates thus represented directly.  It is fun that we are using general set theoretical
objects as variables.  Of course we are only really interested in environments with finite domain and predicates with finite Domain, but before we have the power of set comprehension, we cannot naturally even define what this means.  So we do not restrict ourselves.

It may also be the case that we can show under strong assumptions that some infinitary logic works in extensions of NFU.



\end{slide}

\begin{slide}
 
 This can be reconstrued as giving a semantics for formulas of first order logic, with correlated axioms that certain sets always exist.
 
 As initial framework, provide weak extensionality (or strong if desired), empty sets, singletons, complements and binary unions.
 
 Define a variable as a set with more than one element.  We use letters from $p$ to $w$ to stand for variables, by convention.
 
 Define an environment as a collection $e$ of sets $\{\{x\},v\}$ where $v$ is (as our convention suggests) a variable, in which each variable
 is an element of no more than one element of $e$.  We use letters from $d$ to $h$ by convention to represent environments.  Notice that the empty set is an environment.
 
 The domain ${\tt dom}(e)$ of an environment $e$ is the set of all variables which belong to elements of the environment.  We provide by axiom
 that all environments have domains.
 
 We enclose a formula of first order logic in brackets to represent its semantics.  We omit brackets inside brackets.
 
 Define $[u=v]$ as the set of all environments $\{\{\{x\},u\},\{\{x\},v\}\}$.  We provide that this exists for each $u,v$.
 
 Define $[u,E,v]$ as the set of all environments $\{\{\{\{x\}\},u\},\{\{y\},v\}\}$ where $x \in y$.  This is the semantics for the relation of atomic inclusion.  We provide that this exists for each $u,v$.
 
 Define $[v=a]$ (where $a$ is any constant) as $\{\{\{a\},v\}\}$.  This doesn't need an axiom.
 
 Define a predicate as a set of environments each element of which has the same domain.  We say that $\Delta$ is a Domain of a predicate $P$ just in case $\Delta$ is the domain of each element of $P$, defining the Domain of the empty set as the empty set.    Notice that atomic formulas defined just above are predicates.
 
 Where $P$ is a predicate, define $[\neg P]$ as the set of all environments with domain the Domain of $P$ which do not belong to $P$.
 Notice that $[\neg[\neg P]]=P$ unless $P$ contains all environments with domain its Domain, in which case $[\neg[\neg P]]=\{\emptyset\}$
Note that by convention we can write just $[\neg\neg P]$. We provide by axiom that negations of predicates exist.

Where $P$ and $Q$ are predicates, define $[P \wedge Q]$ as the set of all environments of the form $e \cup f$ where $e \in P$ and $f \in Q$.  Notice that not all sets of the form $e \cup f$ with $e \in P$ and $f \in Q$ are necessarily environments:  the requirement that
elements of this set be environments is a genuine restriction.  We provide by axiom that conjunctions of predicates exist.

We can then define semantics for all binary propositional connectives using the usual definitions in terms of negation and conjunction.

Where $P$ is a predicate and $v$ is a variable, define $[(\exists v:P)]$ as the set of all $$e \setminus \{\{\{x\},v\}:x=x\}$$ such that $e \in P$.  We provide by axiom that existential quantifications of predicates exist,
and note that universal quantifications of predicates can then be defined in a standard way.

We provide for each environment $e$ the environment $e^\iota = \{\{\{\{x\}\},\{v,\emptyset\}\}:\{\{x\},v\}\in P\}$.  We provide for
each predicate $P$ the predicate $P^\iota = \{e^\iota:e \in P\}$.

We provide for each relation $P$ and variable $v$ the set $\{x:\{\{x\},v\}\in P\}$  (the use of this is to extract the extensions of predicates with singleton Domain, but there is no reason to restrict the statement of the axiom).

We provide the axiom of set union.

Observe that by construction the semantics $[\phi]$ of a formula for which we have provided semantics is the set of all environments representing assignments of values to its free variables which make it true, possibly omitting variables which can be assigned any value desired.

This does not immediately give us semantics for stratified formulas, because we do not have semantics for membership.  The key is to convert any stratified formula to a formula with the same truth value for which we do have semantics.  If $\phi$ is stratified with stratification $\sigma$ with the maximum of its range equal to $N$,
replace each component membership formula $u \in v$ with $[u\,E\,v]^{\iota^{N-\sigma(v)}}$, replace each formula $u=v$ with
$[u=v]^{\iota^{N-\sigma(v)}}$, and convert each quantified formula $(\exists v:\psi)$ to $$(\exists v:[v=v]^{\iota^{N-\sigma(v)}}\wedge \psi^*)$$ ($\psi^*$ the already transformed form of $\psi$; with universal quantification considered as defined in terms of existential quantification).  This has the effect of changing all references to values $x$  assigned 
to variables $v$ to references to $\iota^{N-\sigma(v)}(x)$ [stratification of the formula ensures that the same number of $\iota$'s are affixed to each occurrence of a value assigned to a variable -- the variables themselves are actually changed by the same number of iterated applications of $v \mapsto \{v,\emptyset\}$ to each occurrence, but this can be neglected], without changing the truth value of the formula represented.  A unary stratified predicate $\phi(v)$ with stratification $\sigma$ transforms to a predicate of $\iota^{N-\sigma(v)}(v)$ with semantics (speaking a little loosely), and applying set union $N-\sigma(v)$ times to the extension of this predicate gives
$\{x:\phi(x)\}$.

Semantics for constants allow us to handle parameters.

\end{slide}

\end{document}

